% TOP_PROBLEM: {"id": 3217, "title": "Jaeger's modular orientation conjecture", "category_id": 3, "category": {"id": 3, "name": "graph_theory", "display_name": "Graph Theory", "description": "Problems involving graphs, networks, and their properties.", "slug": "graph-theory", "order_index": 3, "created_at": "2026-07-31T15:26:25.671Z"}, "status": "open", "problem_number": "OPG-130", "set_id": 12}
% FIRST_POSTED: 2026-10-01
% ATTEMPT_STATUS: unresolved
% MODEL: GPT 6 Astra Ultra
% UnsolvedMath ID 3217; source catalogue code OPG-130
% !TeX program = lualatex
% Research attempt dated 2026-10-01; canonical statement preserved.
\documentclass[11pt,a4paper]{article}
\usepackage[margin=25mm]{geometry}
\usepackage{fontspec}
\setmainfont{lmroman10-regular.otf}[BoldFont=lmroman10-bold.otf,
 ItalicFont=lmroman10-italic.otf,BoldItalicFont=lmroman10-bolditalic.otf]
\usepackage{amsmath,amssymb,mathtools}
\usepackage{unicode-math}
\setmathfont{latinmodern-math.otf}
\AtBeginDocument{\let\setminus\smallsetminus}
\usepackage{xurl,hyperref}
\hypersetup{colorlinks=true,urlcolor=blue}
\setcounter{secnumdepth}{0}
\setlength{\parindent}{0pt}
\setlength{\parskip}{5pt}
\setlength{\emergencystretch}{3em}
\begin{document}
\section*{Jaeger's modular orientation conjecture}\label{3217}
{\small UnsolvedMath ID 3217 · OPG-130 · Graph Theory · OpenGarden}
\subsection{Definitions and mathematical statement}
Let $k\geq1$ be an integer and
$G=(V,E)$ a finite undirected loopless
multigraph with at least two vertices.
For $\varnothing\neq X\subsetneq V$,
let $\delta_G(X)$ denote the set of
edges with exactly one endpoint in
$X$, counting parallel edges separately.
Assume that $G$ is $4k$-edge-connected:
\[
 |\delta_G(X)|\geq4k
 \qquad(\varnothing\neq X\subsetneq V).
\]

An orientation assigns one direction
to each edge. Write $d^+(v)$ for the
number directed out of $v$ and
$d^-(v)$ for the number directed
into $v$. A modular $(2k+1)$-orientation
is an orientation for which
\[
 d^-(v)-d^+(v)\equiv0\pmod{2k+1}
 \qquad(v\in V).
\]
The difference need not vanish as
an integer; it must be divisible
by the indicated odd modulus.

Jaeger's conjecture asserts that
every graph satisfying the cut
condition has a modular
$(2k+1)$-orientation, for every
integer $k\geq1$. All edge
directions must be chosen together
to meet every vertex congruence.
There is no choice of variable
edge weights in this formulation:
each directed edge contributes
one at its head and minus one
at its tail.

Equivalently, with these fixed
unit values the oriented incidence
equations vanish in the cyclic
group $\mathbb Z/(2k+1)\mathbb Z$.
The edge-connectivity hypothesis
constrains every cut, rather than
only vertex degrees. Neither
regularity nor planarity is assumed.
\subsection{Short English statement}
Can every 4k-edge-connected graph be oriented so each vertex has incoming and outgoing degrees congruent modulo 2k+1?
\subsection{Research attempt}
\paragraph{Scope and method.}
This research attempt by GPT-6 Astra Ultra is dated 1 October 2026.
It preserves the catalogue statement and distinguishes elementary proofs
below from cited prior work. No novelty claim or formal proof-assistant
verification is made. The final paragraph states the precise remaining scope.
\paragraph{Reformulating the exact demand.}
Source [S2] attributes the conjecture to Jaeger. We investigate prescribed
outdegrees rather than replace unit edge contributions by arbitrary
nonzero flow values. Put $m=2k+1$. At a vertex of degree $d_v$, an
outdegree $a_v$ gives imbalance $d_v-2a_v$. Thus the modular condition
is exactly
\[
 a_v\equiv 2^{-1}d_v\pmod m,\qquad 0\le a_v\le d_v.
\]
Here $2^{-1}$ exists because $m$ is odd. Choosing an allowed residue at
each vertex independently does not yet produce an orientation.

\paragraph{Necessary cut arithmetic.}
For $X\subseteq V$, internal edges cancel when vertex imbalances are
summed. If $c=|\delta(X)|$ and $b$ cut edges point out of $X$, the sum
is $c-2b$. It must be divisible by $m$. In particular an odd cut of
size less than $m$ is impossible: its imbalance is odd with absolute
value below $m$, and therefore cannot be a multiple of $m$.
The proposed $4k=2m-2$ edge-connectivity threshold excludes these small
cuts, but the absence of this arithmetic obstruction does not itself
make all cut requirements compatible.

\paragraph{A complete orientation criterion for fixed outdegrees.}
For integer candidates $a_v$, an orientation with those outdegrees exists
if and only if
\[
 \sum_{v\in V}a_v=|E|,\qquad
 \sum_{v\in X}a_v\ge |E(G[X])|\quad(X\subseteq V),
 \qquad a_v\ge0.
\]
Necessity counts the tails of internal edges. To prove sufficiency, create
one left node for each edge and $a_v$ right slots at each vertex $v$.
Join an edge node to all slots at its two endpoints. For any set $F$ of
edge nodes, let $X$ be its endpoint set. It has $\sum_{v\in X}a_v$
neighbouring slots, at least $|E(G[X])|\ge |F|$. Hall's matching criterion
therefore matches every edge to one slot. Orient that edge away from the
vertex owning the slot. The total number of slots is $|E|$, so every
slot is used and all prescribed outdegrees are attained. Parallel edges
are separate left nodes and cause no change in the proof.

The conjecture is consequently reduced to selecting the residues' integer
lifts $a_v=r_v+mt_v$ so that all these inequalities hold. The criterion
gives an exact certificate and a finite matching algorithm once the lifts
are supplied. The unresolved part is choosing them under the stated
edge-connectivity hypothesis, not solving the final matching instance.

\paragraph{Two classes and a sharp warning at small connectivity.}
If every degree is even, an Euler orientation has equal in- and outdegree
at every vertex and works for every $k$, irrespective of connectivity.
For two vertices joined by $N$ parallel edges, take $a=N/2$ edges in one
direction when $N$ is even. If $N$ is odd and $N\ge m$, take
$a=(N-m)/2$. The two imbalances are then $m$ and $-m$. Hence every
two-vertex instance with $N\ge4k$ satisfies the conjecture.

At $k=1$, the graph $K_4$ shows why arbitrary cubic graphs cannot replace
four-edge-connected graphs. Degree three and modulus three force every
vertex to be either a source or a sink. Every edge must join opposite
types, which would give a bipartition of $K_4$, an impossibility.
Its edge-connectivity is three, below the hypothesis, so this is a
threshold warning rather than a counterexample.

\paragraph{Verification and final status.}
The root batch script enumerates the $2^6$ orientations of $K_4$ to check
the obstruction and verifies the two-vertex formulas for $1\le k\le10$
and $4k\le N\le4k+25$. No numerical test replaces the universal cut
criterion. \textbf{UNRESOLVED.} The orientation problem has an exact
integer-lift and matching formulation, and two families are proved.
We do not establish feasible lifts simultaneously for every subset in
an arbitrary $4k$-edge-connected graph.

\subsection{Sources}
\begin{itemize}
\item[S1] ULAM.AI and the UnsolvedMath Contributors, supplied problems.json, record 3217 (OPG-130), OpenGarden collection. Catalogue identity and source transcription; CC BY 4.0.
\url{https://huggingface.co/datasets/ulamai/UnsolvedMath}.
\item[S2] Open Problem Garden, Jaeger's modular orientation conjecture. Original problem and discussion; the mathematical formulation below is expanded from the supplied source record.
\url{http://www.openproblemgarden.org/op/jaegers_modular_orientation_conjecture}.

\end{itemize}
\end{document}
